% Preámbulo modular para presentaciones Beamer (Modelación Estadística)
% Diseño y paleta institucional del Tecnológico de Monterrey

\documentclass[aspectratio=169,xcolor={dvipsnames,table}]{beamer}

\usepackage[utf8]{inputenc}
\usepackage[spanish,mexico]{babel}
\usepackage{amsmath,amssymb,amsfonts,mathtools}
\usepackage{graphicx}
\usepackage{listings}
\usepackage{booktabs}
\usepackage{tikz}
\usetikzlibrary{matrix,arrows,decorations.pathmorphing,shapes.geometric,positioning}

% Paleta Institucional Tecnológico de Monterrey
\definecolor{TecAzul}{HTML}{0039A6}       % Azul TEC: color primario institucional
\definecolor{TecAzulOscuro}{HTML}{1A2E51} % Azul marino oscuro
\definecolor{TecRojo}{HTML}{EC2661}       % Rojo de acento (paleta miTec)
\definecolor{TecGris}{HTML}{646464}       % Gris neutro para comentarios y subtítulos
\definecolor{TecFondoBloque}{HTML}{F4F6F9}% Fondo suave para bloques

% Configuración de Tema Beamer
\usetheme{Boadilla}
\usecolortheme[named=TecAzulOscuro]{structure}
\setbeamercolor{alerted text}{fg=TecRojo}
\setbeamercolor{palette primary}{bg=TecAzulOscuro,fg=white}
\setbeamercolor{palette secondary}{bg=TecAzul,fg=white}
\setbeamercolor{palette tertiary}{bg=TecAzulOscuro!80!black,fg=white}
\setbeamercolor{title}{fg=TecAzulOscuro,bg=TecFondoBloque}
\setbeamercolor{frametitle}{fg=TecAzulOscuro,bg=TecFondoBloque}

% Bloques personalizados elegantes
\setbeamercolor{block title}{bg=TecAzulOscuro,fg=white}
\setbeamercolor{block body}{bg=TecFondoBloque,fg=black}
\setbeamercolor{block title alerted}{bg=TecRojo,fg=white}
\setbeamercolor{block body alerted}{bg=TecRojo!8,fg=black}
\setbeamercolor{block title example}{bg=TecAzul,fg=white}
\setbeamercolor{block body example}{bg=TecAzul!8,fg=black}

% Redondear bordes y añadir sombra sutil
\setbeamertemplate{blocks}[rounded][shadow=true]
\setbeamertemplate{navigation symbols}{}

% Pie de página limpio con número de diapositiva
\setbeamertemplate{footline}{
  \leavevmode%
  \hbox{%
  \begin{beamercolorbox}[wd=.7\paperwidth,ht=2.25ex,dp=1ex,left]{author in head/foot}%
    \hspace*{2ex}\insertshortauthor~~(\insertshortinstitute) \hfill \insertshorttitle
  \end{beamercolorbox}%
  \begin{beamercolorbox}[wd=.3\paperwidth,ht=2.25ex,dp=1ex,right]{date in head/foot}%
    \insertshortdate{}\hspace*{2em}
    \insertframenumber{} / \inserttotalframenumber\hspace*{2ex}
  \end{beamercolorbox}}%
  \vskip0pt%
}

% Configuración de Listings de Python para visualización pedagógica de código
\lstset{
    language=Python,
    basicstyle=\ttfamily\footnotesize,
    keywordstyle=\color{TecAzulOscuro}\bfseries,
    stringstyle=\color{TecRojo},
    commentstyle=\color{TecGris}\itshape,
    showstringspaces=false,
    breaklines=true,
    frame=lines,
    rulecolor=\color{TecAzulOscuro!30},
    backgroundcolor=\color{TecFondoBloque},
    aboveskip=0.5em,
    belowskip=0.5em,
    literate={á}{{\'a}}1 {é}{{\'e}}1 {í}{{\'i}}1 {ó}{{\'o}}1 {ú}{{\'u}}1
             {Á}{{\'A}}1 {É}{{\'E}}1 {Í}{{\'I}}1 {Ó}{{\'O}}1 {Ú}{{\'U}}1
             {ñ}{{\~n}}1 {Ñ}{{\~N}}1 {¿}{{?`}}1 {¡}{{!`}}1
}

% Metadatos institucionales por defecto
\title[Modelación Estadística]{Modelación Estadística para Ciencia de Datos e Ingeniería}
\author[J. Castillo Colmenares]{Juliho David Castillo Colmenares}
\institute[Tec de Monterrey]{Tecnológico de Monterrey \\ Escuela de Ingeniería y Ciencias}
\date{\today}

% Comandos y macros estadísticas compartidas para las presentaciones Beamer (Metropolis)

% Conjuntos numéricos básicos
\newcommand{\R}{\mathbb{R}}
\newcommand{\N}{\mathbb{N}}
\newcommand{\Q}{\mathbb{Q}}
\newcommand{\Z}{\mathbb{Z}}
\newcommand{\set}[1]{\left\{ #1 \right\}}
\newcommand{\abs}[1]{\left|#1\right|}
\newcommand{\norm}[1]{\left\| #1 \right\|}

% Comandos estadísticos y probabilísticos del curso
\newcommand{\Var}{\operatorname{Var}}
\newcommand{\cov}{\operatorname{Cov}}
\newcommand{\corr}{\rho}
\newcommand{\comb}[2]{\begin{pmatrix} #1 \\ #2 \end{pmatrix}}
\newcommand{\s}{\sigma}
\newcommand{\particion}{\mathfrak{P}}
\newcommand{\card}[1]{\#\left( #1 \right)}
\newcommand{\seq}[3]{\set{#1}_{#2}^{#3}}
\newcommand{\E}{\mathbb{E}}
\newcommand{\Prob}{\mathbb{P}}

% Entornos pedagógicos y cajas en español académico natural
\newenvironment{discusion}{%
  \begin{alertblock}{Pregunta de Discusión}%
}{%
  \end{alertblock}%
}

\newenvironment{reflexion}{%
  \begin{alertblock}{Reflexión para el Aula}%
}{%
  \end{alertblock}%
}

\newenvironment{paradoja}{%
  \begin{alertblock}{Paradoja de Exploración}%
}{%
  \end{alertblock}%
}

\newenvironment{motivacion}{%
  \begin{exampleblock}{Motivación: Ciencia de Datos en la Práctica}%
}{%
  \end{exampleblock}%
}

\newenvironment{retoclase}{%
  \begin{block}{Reto Rápido en Equipo}%
}{%
  \end{block}%
}


\title[Introducción a la Probabilidad]{Introducción a la Probabilidad}
\subtitle{Modelación Matemática del Azar y la Incertidumbre}
\author[J. Castillo Colmenares]{Juliho Castillo Colmenares}
\institute{Tecnológico de Monterrey}
\date{}

\begin{document}

% Slide 1: Portada
\begin{frame}[plain]
  \titlepage
\end{frame}

% Slide 2: Hoja de Ruta e Indicador de Posición
\begin{frame}{Hoja de Ruta — Capítulo 02: Teoría de Probabilidad}
  ¿Dónde estamos en nuestro recorrido por la teoría de probabilidad? \pause
  \begin{itemize}
    \item \color{TecRojo}\textbf{02.00 Introducción a la probabilidad \emph{(Hoy)}}\color{black} \pause
    \item \textbf{02.01 Teoría de conjuntos y probabilidad} \emph{(Siguiente sesión)} \pause
    \item \textbf{02.02 Fundamentos de probabilidad} \pause
    \item \textbf{02.03 Técnicas de conteo} \pause
    \item \textbf{02.04 Probabilidad condicional} \pause
    \item \textbf{02.05 Teorema de Bayes} \pause
    \item \textbf{02.06 Muestreo aleatorio}
  \end{itemize}
  \vspace{0.8em}
  \begin{block}{Objetivo de la Sección 2.1}
    Comprender la necesidad de modelar fenómenos aleatorios donde el determinismo clásico falla y establecer la conexión deductiva/inductiva entre Probabilidad y Ciencia de Datos.
  \end{block}
\end{frame}

% Slide 3: Motivación / Dilema Clásico (Pregunta Gatillo)
\begin{frame}{¿Por qué no podemos predecir todo con exactitud?}
  \begin{motivacion}
    Si conocemos la gravedad, la masa de un cohete y el empuje de sus motores, la física clásica nos permite predecir su trayectoria con exactitud milimétrica. Sin embargo, ¿podemos predecir el minuto exacto en que un usuario de \textbf{Amazon} hará clic en "Comprar"? ¿O si un paciente responderá positivamente a un nuevo tratamiento oncológico?
  \end{motivacion}
  
  \vspace{0.6em}
  \pause
  \begin{itemize}
    \item En sistemas biológicos, sociales y tecnológicos complejos, el número de variables incontrolables o no observables es gigantesco. \pause
    \item La \textbf{Teoría de Probabilidad} nos brinda el lenguaje matemático exacto para gestionar y razonar bajo esta incertidumbre ineludible.
  \end{itemize}
\end{frame}

% Slide 4: Fenómenos Deterministas vs Aleatorios
\begin{frame}{Fenómenos Deterministas vs. Fenómenos Aleatorios}
  \begin{columns}[T]
    \begin{column}{0.48\textwidth}
      {\large \color{TecAzulOscuro}\textbf{Fenómeno Determinista}}
      \vspace{0.3em}
      \begin{itemize}
        \item \textbf{Definición:} Bajo las mismas condiciones iniciales, el resultado es siempre idéntico y predecible con certeza. \pause
        \item \textbf{Ejemplos:} La caída de un objeto en el vacío ($d = \frac{1}{2}gt^2$), el área de un círculo ($\pi r^2$), el interés compuesto a tasa fija. \pause
        \item \textbf{Incertidumbre:} $0\%$.
      \end{itemize}
    \end{column}
    
    \begin{column}{0.48\textwidth}
      {\large \color{TecRojo}\textbf{Fenómeno Aleatorio}}
      \vspace{0.3em}
      \begin{itemize}
        \item \textbf{Definición:} Su resultado no puede predecirse de antemano con certeza, incluso reproduciendo las condiciones iniciales. \pause
        \item \textbf{Ejemplos:} El tiempo de latencia en un servidor web, el precio de una acción, el lanzamiento de un dado. \pause
        \item \textbf{Incertidumbre:} Estructurada por regularidades estadísticas a largo plazo.
      \end{itemize}
    \end{column}
  \end{columns}
\end{frame}

% Slide 5: La Relación Probabilidad vs Estadística
\begin{frame}{El Puente Dual: Probabilidad vs. Estadística}
  La probabilidad y la estadística son las dos caras de la misma moneda del razonamiento científico:
  
  \vspace{0.2em}
  \begin{block}{El Motor Deductivo vs. El Telescopio Inductivo}
    \begin{itemize}
      \item \textbf{Probabilidad (Deducción):} Parte de un \emph{modelo teórico poblacional} conocido y predice la verosimilitud (\emph{likelihood}) de los datos o muestras que podríamos observar. \pause
      \item \textbf{Estadística Inferencial (Inducción):} Parte de los \emph{datos muestrales observados} y utiliza las leyes de probabilidad para inferir o estimar el modelo poblacional subyacente.
    \end{itemize}
  \end{block}
  
  \vspace{0.2em}
  \pause
  \begin{center}
    \small \textbf{Sin teoría de probabilidad, la estadística inferencial carecería de cimientos rigurosos para medir el margen de error.}
  \end{center}
\end{frame}

% Slide 6: Campos de Aplicación en Ciencia e Ingeniería
\begin{frame}{Áreas de Aplicación en la Era de los Datos}
  La probabilidad es el núcleo de las tecnologías más avanzadas:
  \vspace{0.4em}
  \begin{itemize}
    \item \textbf{Inteligencia Artificial y Machine Learning:} Redes neuronales probabilísticas, algoritmos de clasificación de Naive Bayes y aprendizaje por refuerzo. \pause
    \item \textbf{Ingeniería de Confiabilidad (SRE):} Modelado del tiempo entre fallas de servidores críticos, tolerancia a fallos y teoría de colas. \pause
    \item \textbf{Finanzas Cuantitativas:} Valoración de opciones de Black-Scholes, Valor en Riesgo ($VaR$) y modelado de portafolios de inversión. \pause
    \item \textbf{Bioinformática y Genética:} Decodificación del genoma humano, mutaciones evolutivas y ensayos clínicos doble ciego.
  \end{itemize}
\end{frame}

% Slide 7: Un Poco de Historia y Contexto
\begin{frame}{Evolución Histórica de la Probabilidad}
  \begin{itemize}
    \item \textbf{Siglo XVII (El Nacimiento):} En $1654$, una correspondencia epistolar entre \textbf{Blaise Pascal} y \textbf{Pierre de Fermat} sobre la justa repartición de apuestas interrumpidas en juegos de azar sentó las bases algebraicas de la probabilidad matemática. \pause
    \item \textbf{Siglos XVIII y XIX (Consolidación):} \textbf{Jacob Bernoulli} demuestra la primera Ley de los Grandes Números ($1713$) y \textbf{Carl Friedrich Gauss} caracteriza la distribución normal del error de medición ($1809$). \pause
    \item \textbf{Siglo XX (Axiomatización Rigurosa):} En $1933$, \textbf{Andrey Kolmogorov} publica \emph{Los Fundamentos de la Teoría de la Probabilidad}, unificando la probabilidad con la teoría de la medida de Lebesgue en solo tres axiomas matemáticos incontestables.
  \end{itemize}
\end{frame}

% Slide 8: Estructura Conceptual del Capítulo 02
\begin{frame}{Estructura del Capítulo: Lo que Construiremos}
  En las siguientes secciones dominaremos las herramientas formales del modelado aleatorio:
  \vspace{0.6em}
  \begin{enumerate}
    \item \textbf{Álgebra de Eventos (Sección 2.2):} Uso de conjuntos, diagramas de Venn y particiones para dividir y contar el universo de posibilidades sin errores. \pause
    \item \textbf{Axiomas de Kolmogorov (Sección 2.3):} Las tres reglas fundamentales de las cuales se deriva toda fórmula probabilística. \pause
    \item \textbf{Probabilidad Condicional (Sección 2.4):} Cómo ajustar la probabilidad ante la presencia de información parcial u observaciones previas. \pause
    \item \textbf{Teorema de Bayes (Sección 2.5):} El motor del aprendizaje automático que actualiza creencias a la luz de nueva evidencia.
  \end{enumerate}
\end{frame}

% Slide 9: Resumen y Conexión
\begin{frame}{Resumen y Conexión con la Sección 2.2}
  \begin{block}{Lo que debemos recordar}
    \begin{itemize}
      \item Los fenómenos aleatorios presentan variabilidad inherente cuya estructura a largo plazo se describe mediante leyes probabilísticas. \pause
      \item La probabilidad es el cimiento matemático de la estadística inferencial y de todo algoritmo moderno de Machine Learning.
    \end{itemize}
  \end{block}
  
  \vspace{0.8em}
  \pause
  {\large \color{TecAzulOscuro}\textbf{Siguiente Sección: 02.02 Notación de Conjuntos y Particiones}}
  \vspace{0.3em}
  \begin{itemize}
    \item Para calcular probabilidades sin ambigüedades, primero debemos aprender a definir el universo de resultados algebraicamente utilizando \textbf{Conjuntos} y descomponerlo como un rompecabezas con \textbf{Particiones}.
  \end{itemize}
\end{frame}

\end{document}
